\section{Discussion}
\label{sec:discussion}

\subsection{What the negative result means}

Thirteen factors, roughly fifty training runs, and no improvement on a stock
configuration beyond run-to-run variance. It is worth being precise about what
this does and does not establish.

It does not establish that CAMUS segmentation is solved, nor that architecture
is irrelevant in general. It establishes something narrower and, we think, more
useful: on this benchmark, at this data scale, the design choices practitioners
routinely tune --- capacity, resolution, mask-head resolution, label encoding,
augmentation recipe, loss weighting, optimiser --- do not move boundary accuracy
by more than the seed changes it. A practitioner is better served by running
three seeds of a stock configuration than by tuning one.

The factors that did register are informative by contrast. Both are
\emph{training-side} and both concern the information available to the model
rather than its capacity to represent: augmentation ($0.318$\,mm) and COCO
initialisation ($0.183$\,mm). Adding parameters, pixels or prototype resolution
consistently did nothing or hurt. On 1600 images the model is regularisation-
and information-bound, not capacity-bound.

\subsection{Where the remaining error lives}

The gap between the achieved $3.809$\,mm HD95 and the $0.31$\,mm encoding ceiling
is a factor of twelve, and none of our interventions closed any of it. Three
independent observations point at the annotations rather than the model.

First, mean Dice already exceeds the inter-observer agreement reported by the
dataset's authors ($0.899$). Above that level the metric increasingly rewards
conformity to one annotator's habits.

Second, the boundary-weighted loss --- which upweights exactly the pixels where
the contour lies --- was the one intervention that was significantly
\emph{harmful}. If the contour were where recoverable signal remained, weighting
it should have helped. That it hurts is consistent with the contour being where
the reference is least reliable, so that emphasising it amplifies label noise.

Third, the EF analysis reproduces the pattern through an independent clinical
endpoint: perfect segmentation yields $5.44\%$ EF MAE, and the best model
$6.66\%$. Most of the error survives perfect masks.

We were unable to test this directly. CAMUS released a multi-observer subset ---
a second and third cardiologist contoured 50 patients, and the original
annotator repeated the task seven months later --- but that data is not part of
the public distribution. We enumerated the public collection and confirmed it
contains only single-observer ground truth. An implementation that would compute
observer agreement under our exact protocol is included in the released code for
anyone who obtains those contours.

\subsection{Implications for evaluation practice}

Three findings of this study are methodological rather than architectural, and
we believe they matter beyond CAMUS.

\paragraph{Single-seed ablations produce false positives} A paired Wilcoxon test
over the 200-frame CAMUS test set assigns $p{=}0.035$ to two runs that differ
only in random seed. Small test sets and paired tests combine to make seed noise
look significant. Several differences reported in the CAMUS literature are
smaller than the $0.134$\,mm we measure from the seed alone. Reporting a
seed-variance estimate should be a minimum standard on benchmarks of this size.

\paragraph{Evaluation faults can hide in released reference code} The official
CAMUS EF implementation selects a contour by position rather than area. This is
correct on ground truth by construction and therefore invisible to the authors'
own validation, but it reduces EF correlation on model predictions from
$r{=}0.825$ to $r{=}0.080$ at native resolution. A reference implementation
being widely used is not evidence that it behaves correctly on the inputs one is
actually feeding it.

\paragraph{Protocols must be stated, not assumed} We found that computing
boundary metrics on a resampled grid while retaining native spacing
under-reports distances; that volumes computed the same way are physiologically
impossible ($\approx$10\,ml rather than $\approx$108\,ml) yet yield plausible EF
because the error partially cancels in the ratio; and that latency figures in
the literature are not comparable across papers because hardware differs. Each
of these silently changes reported numbers without producing any visible error.

\subsection{Practical recommendations}

For practitioners deploying instance segmentation on cardiac ultrasound:

\begin{itemize}
\item Use the filled-contour decomposition for annular structures. Not because
      it is more accurate --- it is not --- but because it is simpler and its
      ceiling is one pixel.
\item Set the confidence threshold far below the natural-image default. Where
      every structure is known to be present, $0.10$ removed every missed
      detection at no accuracy cost.
\item Do not tune capacity. Spend the compute on seeds instead.
\item Apply a connected-component filter before any volume or EF computation.
\item Measure latency yourself, under one protocol, at the resolution you will
      deploy at.
\end{itemize}

\subsection{Limitations}

This study covers one dataset, one acquisition protocol and one GPU class, with
a single official split rather than cross-validation, so the reported spread
reflects seed rather than split variation. Most ablation arms are single-seed
and are interpreted only against the measured noise floor, never individually;
where an arm approached that floor we have said so rather than claiming an
effect. Test-time augmentation is unavailable for segmentation models in the
framework used --- it is silently ignored --- so it is untested rather than
found ineffective. Latency is measured on a laptop discrete GPU rather than on
the embedded accelerator the deployment argument ultimately points at. Finally,
the annotation-limit hypothesis is inferred from three converging observations
rather than measured directly, because the multi-observer data is not public.
